\section{Linguagens de primeira ordem}\label{sec:linguagensPrimeiraOrdem}
\index{linguagem de primeira ordem}
\index{símbolo!lógico}
\index{símbolo!não lógico}

Discutida a notação polonesa de forma geral, estamos prontos para avançar para a linguagem da lógica de primeira ordem.

Iniciaremos apresentando os símbolos lógicos.
Enfatizamos, desde já, que consideraremos a igualdade como um símbolo lógico, o que é uma escolha comum, mas não universal, na literatura.

\begin{metadefinition}\label{mdef:simbolosLogicos}\index{símbolo!lógico}\index{conectivo lógico}\index{quantificador}
    Os símbolos lógicos (da lógica de primeira ordem) são os seguintes.
    \begin{itemize}
        \item Os símbolos de conectivos lógicos:
        \begin{itemize}
            \item $\neg$ (negação, unário);
            \item $\wedge$ (conjunção, binário);
            \item $\vee$ (disjunção, binário);
            \item $\rightarrow$ (condicional, binário);
            \item $\leftrightarrow$ (bicondicional, binário).
        \end{itemize}
        \item Os símbolos de quantificadores:
        \begin{itemize}
            \item $\forall$ (quantificador universal, binário);
            \item $\exists$ (quantificador existencial, binário).
        \end{itemize}
        \item $=$, o símbolo de igualdade (binário).
        \item Uma infinidade (enumerável) de símbolos de variáveis, que denotaremos por $x$, $y$, $z$, $x_1$, $y_1$, $z_1$, $x_2$, $y_2$, $z_2$, $\dots$ (variáveis, $0$-árias).
        \end{itemize}
\end{metadefinition}

Neste texto, seremos extremamente permissivos com a notação utilizada para os símbolos de variáveis.
Além das letras utilizadas, poderemos utilizar outros símbolos, como $a$, $\beta_3$, $\alpha_5$, etc., ou até mesmo $f$ (comum em análise real: $\forall f\,(f \text{ é função } \rightarrow\dots)$).
A identificação de um símbolo como variável se dará pelo contexto.

Observamos desde já que, neste formalismo, os quantificadores, $\forall$ e $\exists$, são símbolos binários.
De fato, requerem dois objetos para serem utilizados: um símbolo de variável e uma fórmula.

Dito isso, há duas possíveis dúvidas que podem surgir a respeito da Metadefinição~\ref{mdef:simbolosLogicos}.
A primeira diz respeito ao fato de que uma teoria de primeira ordem não possui apenas esses símbolos.
Por exemplo, a Teoria dos Conjuntos necessita do símbolo de pertencimento, $\in$, e a Teoria dos Grupos necessita de um símbolo de operação binária, $\cdot$.
Tais símbolos são chamados de \emph{símbolos não lógicos} e, conforme veremos, compõem o léxico de uma teoria de primeira ordem.

A segunda dúvida diz respeito ao fato de que nem toda expressão na notação polonesa é o que consideramos ser uma fórmula matemática válida.
Por exemplo, considere a expressão na notação polonesa $\wedge x_1 = x_2 \neg x_3$ -- o que, na notação usual, se traduziria como $x_1 \wedge (x_2 = \neg x_3)$ -- não costuma ser aceito como fórmula válida.
De fato, conforme veremos, nem toda expressão na notação polonesa é uma fórmula bem formada de primeira ordem.

Iniciaremos trabalhando a primeira dúvida.

\begin{metadefinition}\label{mdef:lexicoPrimeiraOrdem}\index{léxico!de primeira ordem}\index{símbolo!não lógico}\index{símbolo!não lógico!de predicado}\index{símbolo!não lógico!funcional}
    Um léxico $\mathcal L$ para uma teoria de primeira ordem consiste no seguinte.

    \begin{itemize}
        \item Em uma coleção $\mathcal P$ de símbolos de predicado (ou relacionais), cada um com uma aridade associada.
        \item Em uma coleção $\mathcal F$ de símbolos funcionais, cada um com uma aridade associada.
    \end{itemize}
    As duas coleções acima não podem possuir símbolos em comum, nem possuir símbolos lógicos.
    Os símbolos de $\mathcal P$ e $\mathcal F$ são chamados de símbolos não lógicos de $\mathcal L$.
    Juntamente com os símbolos lógicos da Metadefinição~\ref{mdef:simbolosLogicos}, eles constituem o léxico total da linguagem.
\end{metadefinition}

Exemplos de símbolos de predicado não lógicos incluem $\in$ (pertencimento, comum em teoria dos conjuntos), $<$ (ordem, comum no estudo dos números reais), $\triangle$ (não colinearidade entre três pontos, comum em geometria plana).

Exemplos de símbolos funcionais não lógicos incluem $\cdot$ (operação binária, comum em teoria dos grupos), $+$ (operação binária, comum em teoria dos anéis), $-$ (operação unária, comum em teoria dos anéis).

Dando exemplos explícitos de léxicos:
\begin{example}Consideremos os seguintes léxicos.\label{exa:exemplosLexicos}
    \begin{itemize}
        \item O léxico da teoria dos conjuntos é composto por um símbolo de predicado não lógico binário, $\in$.
        \item O léxico da teoria dos grupos é composto pelo símbolo funcional binário $\cdot$ (operação de grupo), pelo símbolo funcional unário $^{-1}$ (inverso) e pelo símbolo funcional $0$-ário $e$ (identidade).
        \item O léxico da teoria dos anéis é composto por dois símbolos funcionais binários, $+$ (soma) e $\cdot$ (produto), por um símbolo funcional unário, $-$ (oposto aditivo) e por um símbolo de constante, $0$ (identidade aditiva).
    \end{itemize}
\end{example}
Em textos de álgebra, podem existir diferenças na escolha dos símbolos funcionais.
Por exemplo, alguns livros podem dar a entender que a inversão não é um ente primitivo da teoria dos grupos.
Nestes textos, prova-se que em todo grupo, todo elemento possui único inverso, e, então, introduz-se a inversão por definição.
Trataremos de extensões da linguagem via definição mais adiante.

Vale a pena, neste ponto, darmos uma atenção especial aos símbolos de aridade $0$.

Os símbolos funcionais não lógicos de aridade $0$ são chamados de \emph{símbolos de constantes}, ou simplesmente \emph{constantes}, e exemplos deles incluem $\emptyset$ (o conjunto vazio, comumente introduzido via definição em teoria dos conjuntos), $e$ (identidade, comum em teoria dos grupos) e $\pi$ (o número irracional pi, que aparece no estudo dos números reais).
Ressaltamos que, apesar de variáveis e constantes serem símbolos de aridade $0$, e apesar de estarmos sendo permissivos com a notação utilizada para variáveis, esses objetos são distintos (por exemplo, apenas variáveis podem ser quantificadas!).
A distinção entre as notações de variáveis e constantes é feita pelo contexto, e, em caso de ambiguidade, deve ser indicada sempre que relevante.

Já os símbolos de predicado não lógicos de aridade $0$ são chamados de \emph{letras proposicionais}, e são interpretados como ``sentenças atômicas''.
Tais símbolos podem ser utilizados para formalizar abreviações, como $\AC$ (Axioma da Escolha, comum em teoria dos conjuntos. Do inglês, \emph{Axiom of Choice}) e $\CH$ (Hipótese do Contínuo, comum em teoria dos conjuntos. Do inglês, \emph{Continuum Hypothesis}).

Agora devemos trabalhar a segunda dúvida.
Para isso, devemos definir, precisamente, o que são termos e fórmulas.

Iniciaremos pelos termos.

\begin{metadefinition}\label{mdef:termosPrimeiraOrdem}\index{termo}\index{constante}
    Dado um léxico $\mathcal L$ para uma teoria de primeira ordem, os termos de $\mathcal L$ são definidos recursivamente da seguinte forma.
\begin{enumerate}[label=(\roman*)]
    \item Toda variável é um termo.
    \item Todo símbolo de constante é um termo.
    \item Se $f$ é um símbolo funcional de aridade $n>0$ e $t_1,\dots,t_n$ são termos, então $ft_1\dots t_n$ é um termo.
    \item Todos os termos são obtidos a partir de um número finito de aplicações das regras anteriores.
\end{enumerate}

\end{metadefinition}

Vamos dar alguns exemplos.

\begin{example}\label{exa:termosPrimeiraOrdem}Utilizando a notação do Exemplo~\ref{exa:exemplosLexicos}, temos os seguintes termos.
    \begin{itemize}
        \item Na teoria dos grupos, $\cdot x_1 e$ é um termo (na notação usual, $x_1 \cdot e$).
        \item Na teoria dos anéis, $+ x_1 \cdot -x_20$ é um termo (na notação usual, $x_1 + (-x_2 \cdot 0)$).
        \item Considere o léxico básico da teoria de conjuntos.
        Os únicos termos são as variáveis (adiante, veremos que mais termos podem ser introduzidos por definição).
        \item Na teoria de grupos, $\wedge x_1 e$ não é um termo, pois $\wedge$ é um símbolo de conectivo lógico, e não um símbolo funcional.
    \end{itemize}
\end{example}

Agora podemos definir o que são fórmulas.

\begin{metadefinition}\label{mdef:formulasPrimeiraOrdem}\index{fórmula}\index{fórmula!atômica}\index{fórmula!equacional}\index{subfórmula}
    Dado um léxico $\mathcal L$ para uma teoria de primeira ordem, as fórmulas de $\mathcal L$ são definidas recursivamente da seguinte forma.
\begin{enumerate}[label=(\roman*)]
    \item Se $P$ é um símbolo de predicado de aridade $n>0$ e $t_1,\dots,t_n$ são termos, então $Pt_1\dots t_n$ é uma fórmula.
    \item Se $P$ é um símbolo de predicado de aridade $0$, então $P$ é uma fórmula.
    \item Se $t_1$ e $t_2$ são termos, então ${=}t_1t_2$ é uma fórmula.
    \item Se $\varphi$ é uma fórmula, então $\neg\varphi$ é uma fórmula.
    \item Se $\varphi$ e $\psi$ são fórmulas, então ${\wedge}\varphi\psi$, ${\vee}\varphi\psi$, ${\rightarrow}\varphi\psi$ e ${\leftrightarrow}\varphi\psi$ são fórmulas.
    \item Se  $x$ é uma variável e $\varphi$ é uma fórmula, então $\forall x\varphi$ e $\exists x\varphi$ são fórmulas.
    \item Todas as fórmulas são obtidas a partir de um número finito de aplicações das regras anteriores.
\end{enumerate}

Fórmulas dos tipos (i), (ii) e (iii) são chamadas de \emph{fórmulas atômicas}.
Fórmulas do tipo (iii) são chamadas de \emph{fórmulas equacionais}.
Vale a pena destacar que o símbolo de igualdade, $=$, se comporta sintaticamente exatamente como um símbolo de predicado binário.
Tal símbolo já ``vem no pacote'' da lógica de primeira ordem, diferente do que ocorre, por exemplo, com o símbolo $\in$, presente na Teoria dos Conjuntos.

Finalmente, dizemos que uma \emph{subfórmula} de uma fórmula $\varphi$ é uma subexpressão de $\varphi$ que também é uma fórmula.
\end{metadefinition}

Vamos dar alguns exemplos.

\begin{example}\label{exa:formulasPrimeiraOrdem}Utilizando a notação do Exemplo~\ref{exa:exemplosLexicos}, temos as seguintes fórmulas.
    \begin{itemize}
        \item Na teoria dos grupos, $\cdot x_1 e$ não é uma fórmula, pois $\cdot$ é um símbolo funcional, e não um símbolo de predicado.
        Por outro lado, é um termo (usualmente denotado por $x_1 \cdot e$).
        \item Na teoria dos grupos, $= \cdot x_1 e x_2$ é uma fórmula equacional (na notação usual, $x_1 \cdot e = x_2$).
        \item Na teoria de conjuntos, $\in x_1 x_2$ é uma fórmula atômica (na notação usual, $x_1 \in x_2$).
        \item Em qualquer léxico para uma teoria de primeira ordem ${=}x_1 x_2$ é uma fórmula equacional (na notação usual, $x_1 = x_2$).
        \item Na teoria dos conjuntos, ${\wedge} {\in} x_1 x_2 {=} x_2 x_3$ é uma fórmula (na notação usual, $(x_1 \in x_2) \wedge (x_2 = x_3)$).
        \item Na teoria dos conjuntos, ${\in}{\wedge} x_1 x_2 {=} x_2 x_3$ não é uma fórmula, pois a única chance de uma expressão começada por $\in$ ser fórmula é pela regra (i), mas, nesse caso, deveria haver um termo contendo $\wedge$, o que não é possível.
        Na notação usual, a expressão acima se traduziria como $(x_1 \wedge x_2) \in (x_2 = x_3)$.
        \item No léxico da teoria de conjuntos, $\forall x_1 {\in} x_1 x_2$ é uma fórmula (na notação usual, $\forall x_1 (x_1 \in x_2)$).
        \item No léxico da teoria dos grupos, $\exists e {=} e\cdot e$ não é uma fórmula, pois $e$ é um símbolo de constante, e não uma variável, e a única forma de uma fórmula iniciada por $\exists$ ser uma fórmula é pela regra (vi).
        Na notação usual, a expressão acima se traduziria como $\exists e (e = e \cdot e)$.
    \end{itemize}
\end{example}

Agora sabemos, com precisão, o que são fórmulas.

A partir deste ponto, deixaremos de escrever termos e fórmulas na notação polonesa, a menos que isso seja necessário.
Seguiremos as seguintes convenções.

\begin{convention}\label{conv:convencoesEscritaPrimeiraOrdem}Utilizaremos as seguintes convenções de escrita.
    Escreveremos:
\begin{itemize}
    \item $\phi \wedge \psi$ em vez de ${\wedge}\phi\psi$;
    \item $\phi \vee \psi$ em vez de ${\vee}\phi\psi$;
    \item $\phi \rightarrow \psi$ em vez de ${\rightarrow}\phi\psi$;
    \item $\phi \leftrightarrow \psi$ em vez de ${\leftrightarrow}\phi\psi$;
    \item $t_1 = t_2$ em vez de ${=}t_1 t_2$;
    \item $P(t_1, \dots, t_n)$ em vez de $P t_1 \dots t_n$, em que $P$ é um símbolo de predicado de aridade $n>0$ e $t_1, \dots, t_n$ são termos;
    \item $f(t_1, \dots, t_n)$ em vez de $f t_1 \dots t_n$, em que $f$ é um símbolo funcional de aridade $n>0$ e $t_1, \dots, t_n$ são termos.
\end{itemize}
Utilizaremos parênteses para eliminar ambiguidades sempre que for necessário.
\end{convention}

\subsection*{Exercícios}

\begin{exercise}\label{exr:gruposNotacaoPolonesa}
    Considere a linguagem da Teoria dos Grupos do Exemplo~\ref{exa:exemplosLexicos}.
    Escreva as expressões abaixo, que estão escritas na notação polonesa, na notação usual, e diga quais delas são termos, quais são fórmulas, e quais não são nem termos, nem fórmulas.

    \begin{enumerate}[label=(\roman*)]
        \item $\cdot x_1 e$;
        \item $= \cdot x_1 e x_2$;
        \item ${\wedge} {\cdot} x_1 e {=} x_2 x_3$;
        \item $\exists e {=} e{\cdot} e e$;
        \item $\cdot \cdot x_1 e x_2$;
        \item $=\wedge \cdot x_1 e x_2$.
    \end{enumerate}
\end{exercise}

\begin{exercise}\label{exr:aneisNotacaoPolonesa}
    Considere a linguagem da Teoria dos Anéis do Exemplo~\ref{exa:exemplosLexicos}.
    Escreva as expressões abaixo, que estão escritas na notação usual, na notação polonesa.
    Decida quais delas são termos e quais são fórmulas.
    \begin{enumerate}[label=(\roman*)]
        \item $x_1 + (-x_2 \cdot 0)$;
        \item $\exists x_1 (x_1 + x_1 = 0)$;
        \item $x_1 + x_2 = x_3$;
        \item $(x_1+x_2)\cdot x_3$;
        \item $\forall x_1 \exists x_2 (x_1 + x_2 = 0)$;
        \item $\forall x_1 (x_1 = x_1) \wedge \exists x_1 (x_1 = x_1)$;
        \item $\forall x_1 (x_1 = x_1 \wedge \exists x_1 (x_1 = x_1))$.
    \end{enumerate}
\end{exercise}

\begin{exercise}\label{exr:expressaoNaoPodeSerTermoEFormula}
    Mostre que nenhuma sequência de símbolos de um léxico para uma teoria de primeira ordem é, ao mesmo tempo, um termo e uma fórmula.
\end{exercise}

\begin{exercise}\label{exr:subexpressaoDeFormula}
    Prove ou dê um contra-exemplo: toda subexpressão de uma fórmula é uma fórmula.
\end{exercise}
\begin{exercise}\label{exr:subexpressoesDeFormula}
    Seja $\varphi$ uma fórmula.
    Mostre que uma subexpressão de $\varphi$ é uma fórmula se, e somente se, ela é o escopo de um conectivo lógico, de um quantificador, de um símbolo de igualdade, ou de um símbolo de predicado.
\end{exercise}
